%% ma: L10-B15-0097
%% lop: 10
%% chuong: 6
%% bai: 15
%% dang: 10.15.5
%% loai: TN4
%% muc: TH
%% dap_an: "D"
%% nguon: "Ảnh giáo viên gửi 25/11/2026 (lô 2): Bài tập trắc nghiệm – Toán 10 Bài 15 Hàm số (sự biến thiên, đồ thị), Phần I Câu 8"
%% kien_thuc: "Lớp 10 Bài 15 (đồng biến, nghịch biến); hàm số bậc nhất y=ax+b đồng biến khi a>0 (lớp 9)"
%% trang_thai: chinh-thuc
\begin{ex}
Cho hàm số $y=(16-m^2)x+3$. Có tất cả bao nhiêu giá trị nguyên dương của $m$ để hàm số đã cho đồng biến trên $\mathbb R$?
\choice{$7$}{$9$}{$6$}{\True $3$}
\loigiai{Cần $16-m^2>0\Leftrightarrow-4<m<4$. Các giá trị nguyên dương: $m=1;2;3$. Có $3$ giá trị.}
\end{ex}
